\section{Introduction}
\IEEEPARstart{F}{abric} inspection is one of the most labour-intensive steps of textile manufacturing. Human inspectors are slow, tire quickly, and miss a substantial fraction of defects, which has motivated two decades of research on automated visual inspection~\cite{kumar2008survey,carrilho2024survey}. The state of the art is dominated by supervised deep detectors and segmenters---improved two-stage detectors~\cite{wu2020cascaded}, efficient single-stage detectors~\cite{zhou2021edds,zhang2022mobilenet}, YOLO-family models on edge hardware~\cite{machado2025textile,nasim2024realworld}, transformer segmenters with Gabor priors~\cite{xu2023swin}, and open-world formulations~\cite{ye2026owdln}. These models reach high accuracy on the fabrics they were trained on, but they require hundreds to thousands of labelled defect instances per fabric family and a GPU-centred training workflow; the scarcity of labelled defects and the cost of retraining for new materials are repeatedly identified as the main obstacles to deployment~\cite{carrilho2024survey}.

The obstacle is acute on the \emph{shop floor of small and medium-sized garment and technical-textile plants}, where a number of seaming or joining machines process short runs of many different fabrics. Our motivating installation is a set of ultrasonic seaming stations, each controlled by an ESP32 micro-controller and an Android application (FabricScope) and equipped with a USB camera that already identifies the fibre and weave before the seam parameters are chosen. A fabric typically stays on a machine for minutes, defects are rare, and there is neither time nor expertise to label data or retrain a model whenever the material changes. Such a deployment is naturally an Industrial Internet of Things (IIoT) system~\cite{sisinni2018iiot}: many low-cost camera nodes, a shared plant network, and at most one modest edge server. The inspection function must therefore (i) learn a new fabric from a few seconds of defect-free material, (ii) set its own decision threshold without defect examples, (iii) run within the frame budget of an inexpensive station node, and (iv) scale to many stations without saturating the uplink or the edge server.

Unsupervised anomaly detection, which learns only from normal samples, satisfies (i) and is an active direction for fabrics~\cite{zhang2026dmglad,ma2025pnemfd,an2025causality}; on generic benchmarks, training-free methods that reuse ImageNet features~\cite{rippel2021gaussian,defard2021padim,roth2022patchcore} are particularly strong. However, these works report threshold-free metrics with hundreds of training images on a GPU, and are silent on (ii)--(iv). Conversely, edge-intelligence research~\cite{shi2016edge,zhou2019edgeintel,li2020edgeai} and fog-based manufacturing inspection~\cite{li2018fog} study where to run inference, but assume supervised models trained offline. To our knowledge, no prior work addresses few-shot, label-free surface inspection \emph{as a networked IIoT service}.

This paper fills that gap with \emph{FabricInspect-IoT}, an open-source three-tier system (station node, edge server, profile registry over MQTT) and an empirical study of its design. Our contributions are:
\begin{enumerate}
\item A \textbf{few-shot, label-free inspection service} in which a fabric is learned from $k$ defect-free frames and two complementary models are built at once: a 5--6\,kB Gabor--colour Gaussian model that runs on the station, and a ResNet-18 nearest-neighbour model (PatchCore-lite) that runs on the edge. Decision thresholds are derived by leave-one-out (LOO) scoring of the $k$ shots, without any defect example.
\item \textbf{Confidence-gated offloading}: the station decides clear cases locally and offloads only frames whose local score falls in an uncertainty band around its LOO threshold. Because both thresholds are label-free, the gate needs no tuning data.
\item A \textbf{profile-distribution mechanism} based on retained MQTT messages, which makes a newly learned fabric available to every station within milliseconds, and an analytical model of uplink load and edge utilisation that predicts how many stations an edge server can support.
\item A \textbf{reproducible evaluation} on the MVTec~AD textures~\cite{bergmann2019mvtec,bergmann2021mvtec} covering few-shot accuracy ($k=1\ldots$all), label-free thresholding, and robustness, together with a live testbed with a production MQTT broker~\cite{light2017mosquitto} and a trace-driven simulator validated against it, which scales the system to 64 stations on emulated Wi-Fi and constrained uplinks and reports end-to-end latency, deadline misses, uplink traffic, and edge utilisation.
\end{enumerate}

The remainder of the paper is organised as follows. Section~\ref{sec:related} reviews related work. Section~\ref{sec:system} presents the architecture and protocol, and Section~\ref{sec:methods} the detectors, the thresholds, the gate, and the load model. Section~\ref{sec:setup} describes the experimental setup, Section~\ref{sec:results} reports the results, and Section~\ref{sec:discussion} discusses limitations before Section~\ref{sec:conclusion} concludes.
